% Part: first-order-logic
% Chapter: axiomatic-deduction
% Section: proof-theoretic-notions

\documentclass[../../../include/open-logic-section]{subfiles}

\begin{document}

\iftag{FOL}
      {\olfileid{fol}{axd}{ptn}}
      {\olfileid{pl}{axd}{ptn}}

\olsection{Gagasan Teoretis-Bukti}

\begin{explain}
Kaya nalika kita nemtokake sawetara gagasan semantis wigati
(\iftag{FOL}{validitas}{tautologi}, akibat, lan bisa dicukupi), saiki
kita nemtokake \emph{gagasan teoretis-bukti} kang cocog. Gagasan iki
ora ditemtokake kanthi nggunakake panyukupan !!{sentence}s ing
!!{structure}s, nanging kanthi nggunakake !!{derivability} utawa
!!{nonderivability} formula-formula tartamtu. Panemuan yen gagasan
kasebut padha iku panemuan wigati. Kasunyatan manawa padha dadi isi
\emph{teorema kasahihan} lan \emph{teorema kelengkapan}.
\end{explain}

\begin{defn}[!!^{derivability}]
!!^a{formula}~$!A$ \emph{bisa !!{derivable}} saka $\Gamma$, ditulis
$\Gamma \Proves !A$, yen ana !!a{derivation} saka~$\Gamma$ kang
pungkasan ing~$!A$.
\end{defn}

\begin{defn}[Teorema]
!!^a{formula}~$!A$ iku \emph{teorema} yen ana !!a{derivation} saka
himpunan kosong kang ngasilake $!A$. Kita nulis $\Proves !A$ yen $!A$
iku teorema, lan $\Proves/ !A$ yen dudu teorema.
\end{defn}

\begin{defn}[Konsistensi]
Himpunan !!{formula}s $\Gamma$ iku \emph{konsisten} yen lan mung yen
$\Gamma\Proves/ \lfalse$; yen ora, himpunan iku \emph{ora konsisten}.
\end{defn}

\begin{prop}[Refleksivitas]
\ollabel{prop:reflexivity}
Yen $!A \in \Gamma$, mula $\Gamma \Proves !A$.
\end{prop}

\begin{proof}
  !!{formula}~$!A$ dhewe dadi !!a{derivation} saka~$!A$ saka~$\Gamma$.
\end{proof}

\begin{prop}[Monotonisitas]
\ollabel{prop:monotonicity}
Yen $\Gamma \subseteq \Delta$ lan $\Gamma \Proves !A$, mula $\Delta
\Proves !A$.
\end{prop}

\begin{proof}
Saben !!{derivation} saka $!A$ saka $\Gamma$ uga dadi !!a{derivation}
saka $!A$ saka~$\Delta$.
\end{proof}

\begin{prop}[Transitivitas]
\ollabel{prop:transitivity}
Yen $\Gamma \Proves !A$ lan $\{!A\} \cup \Delta \Proves
!B$, mula $\Gamma \cup \Delta \Proves !B$.
\end{prop}

\begin{proof}
  Upamana $\{!A\} \cup \Delta \Proves !B$. Banjur ana
  !!a{derivation} $!B_1$, \dots, $!B_l = !B$ saka~$\{!A\} \cup
  \Delta$. Sawetara langkah ing !!{derivation} kasebut bener amarga
  sawijining aturan nuduhake larik sadurunge~$!B_i = !A$. Miturut
  hipotesis, ana !!a{derivation} saka~$!A$ saka~$\Gamma$, yaiku
  !!a{derivation}~$!A_1$, \dots, $!A_k = !A$ kang saben $!A_i$ dadi
  aksioma, !!a{element} saka~$\Gamma$, utawa bener miturut aturan
  inferensi. Saiki gatekna rerangken
  \[
  !A_1, \dots, !A_k = !A, !B_1, \dots, !B_l = !B.
  \]
  Iki !!{derivation} kang bener saka~$!B$ saka $\Gamma \cup \Delta$
  amarga saben $!B_i = !A$ saiki dibenerake kanthi dhasar kang padha
  karo kang mbenerake~$!A_k = !A$. Cathetan koreksi sumber OLPL-025:
  sumber Inggris ngilangake tandha metavariabel formula saka sebutan
  kapindho lan nyebut kabeh dhasar kabeneran minangka aturan, sanadyan
  ukara sadurunge mbedakake aksioma, unsur premis, lan aturan; wujud
  formula lan dhasar umum kasebut dibalekake ing kene.
\end{proof}

Gatekna yen iki mligine ateges: yen $\Gamma \Proves !A$ lan $!A
\Proves !B$, mula $\Gamma \Proves !B$. Uga ndherek yen $!A_1,
\dots, !A_n \Proves !B$ lan $\Gamma \Proves !A_i$ kanggo saben~$i$,
mula $\Gamma \Proves !B$.

\begin{prop}
\ollabel{prop:incons}
$\Gamma$ ora konsisten yen lan mung yen $\Gamma \Proves !A$ kanggo
saben~$!A$.
\end{prop}

\begin{proof}
Gladhen.
\end{proof}

\tagprob{FOL}
\begin{prob}
Buktekna \olref[fol][axd][ptn]{prop:incons}.
\end{prob}
\tagendprob

\tagprob{notFOL}
\begin{prob}
Buktekna \olref[pl][axd][ptn]{prop:incons}.
\end{prob}
\tagendprob

\begin{prop}[Kekompakan]
  \ollabel{prop:proves-compact}
  \begin{enumerate}
  \item Yen $\Gamma \Proves !A$, mula ana subhimpunan winates
    $\Gamma_0 \subseteq \Gamma$ saengga $\Gamma_0 \Proves !A$.
  \item Yen saben subhimpunan winates saka~$\Gamma$ konsisten,
    mula $\Gamma$ konsisten.
  \end{enumerate}
\end{prop}

\begin{proof}
  \begin{enumerate}
    \item Yen $\Gamma \Proves !A$, mula ana rerangken winates
      !!{formula}s $!A_1$, \dots,~$!A_n$ saengga $!A \ident !A_n$
      lan saben $!A_i$ iku aksioma logis, !!a{element} saka~$\Gamma$,
      utawa ndherek saka !!{formula}s sadurunge miturut salah siji
      aturan inferensi kang diidini. Jupuken $\Gamma_0$ minangka
      $!A_i$ kang kalebu ing~$\Gamma$. Banjur !!{derivation} kasebut
      uga dadi !!a{derivation} saka~$\Gamma_0$, mula $\Gamma_0
      \Proves !A$. Cathetan koreksi sumber OLPL-026: bukti Inggris
      mung nyebut modus ponens ing alternatif pungkasan, sanadyan versi
      FOL bab iki uga ngidini aturan kuantor; sebutan umum marang aturan
      inferensi kang diidini njaga bukti kanggo loro tag.
    \item Iki kontraposisi saka~(1) kanggo kasus mligi $!A
      \ident \lfalse$.
\end{enumerate}
\end{proof}

\end{document}
